% Appendix A -- parameter measurement, moved out of section 3 (2026-09-26) to
% keep the body within the MLSys page limit. Nothing here is needed to follow
% the argument; all of it is needed to check it.

\section{Measured parameters and sensitivity}
\label{app:params}

A100-80 PCIe; two models; 41 single-shot configurations plus a
13-configuration $n{=}3$ repeat sweep; distinct random contexts; prefix caching
disabled.
% src: paper/model/p4_llama8b_a100_rep3.json @ 2bfc6ca5e

\paragraph{Weight read is efficient and stable within a model.} $\effw \approx
0.75$ of the measured $1935$\,GB/s peak ($\approx$1450\,GB/s achieved). The
$B{=}1$ steps must have their KV term subtracted before $\effw$ is read off
them: it is 1.6\% of the step at 2K context and 6.0\% at 8K, and leaving it in
manufactures a drift (0.729 falling to 0.694 for Llama-3.1-8B) that is the KV
term growing rather than the weight read degrading. Subtracted, $\effw$ is
0.741/0.736/0.738 for Llama-3.1-8B across 2K--8K and 0.778/0.777 for
Qwen2.5-32B --- flat within each model, and 5.3\% apart between them. We carry
the single value 0.75, which is between them, and \S\ref{sec:model:cancel}
reports what each choice predicts.

\paragraph{Paged KV read is neither.} $\effkv$ degrades by $1.6\times$ with the
KV working set (Table~\ref{tab:effkv}):

\begin{table*}[t]
\centering
\small
\begin{tabular}{lrrrrr}
\toprule
KV working set & 4.3\,GB & 8.6\,GB & 17--34\,GB & 43\,GB & 51.5\,GB \\
\midrule
$\effkv$ & 0.77 & 0.72 & 0.54--0.56 & 0.49 & 0.46 \\
\bottomrule
\end{tabular}
\caption{Paged KV read efficiency is not a constant. Error bars from $n{=}3$
repeats are $\pm0.02$--$0.04$; the degradation is far outside them. The
17--34\,GB column takes 0.55 rather than the 0.747 measured at
$B{=}16,S{=}8192$ --- the anomaly discussed below. That is the
\emph{optimistic} choice for us, since a lower $\effkv$ lengthens the step and
so raises the projected gain, and we note it rather than let it pass.}
\label{tab:effkv}
\end{table*}

This is the real correction to the model. It also bears on how such models are
built elsewhere: \textbf{a roofline that charges KV at one fixed bandwidth will
be optimistic at the capacity-bound operating point}, because that is where
$\effkv$ is worst. We measured this on one GPU and one kernel generation
(A100, FlashAttention-2 era), so we state it as a property of what we measured
rather than of the literature.

\paragraph{A caveat at kernel granularity.} The claim that only total KV bytes
matter holds at three of four same-byte test points but breaks at 17.2\,GB
($B{=}32,S{=}4096$ gives 0.561 against $B{=}16,S{=}8192$ at 0.747). The cause
is a reproducible anomaly at $B{=}32,S{=}4096$ that also appeared in an
independent earlier run --- an attention-kernel scheduling artifact, most
likely a split-KV heuristic in the attention kernel choosing a different
partitioning at the larger batch, not measurement noise ($\pm0.61$\,ms
against a $\sim$3\,ms deviation). The aggregate law holds; the
per-configuration surface has kernel-regime steps.

\paragraph{Compute is not the binding term.} Decode-attention arithmetic
intensity equals the GQA group size --- 8 for 70B, 4 for 8B --- against an
H100 balance point near 116 FLOP/byte, so attention sits 14--29$\times$ below
the compute roof. At batch 8, 70B spends 7\% of its step on arithmetic.
% src: paper/model/compute_check.py @ 2bfc6ca5e
Llama-8B at 4K crosses the arithmetic-intensity balance point at batch 142, while its step remains memory-bound by about $4\times$, which is a further
reason it is the wrong model to evaluate on.

